
\subsection{Results: mechanism confirmed, success not}\label{sec:stlsres}

Table~\ref{tab:stlsres} reports the comparison on the straightened models with all four methods
run on identical episodes. Every pre-registered mechanism observable moved, and in the predicted
direction: gradient steps per adaptation update fell from $2.32$ to $1.11$ on the shape family
(predicted $1.2$--$1.5$) and from $2.10$ to $1.32$ on the dynamics family; the effective step size
stopped sitting at the saturating $1.6\times10^{-3}$ and dropped to $2\times10^{-4}$; refinement
depth fell from $2.86$ to $1.06$, i.e.\ recursion is now spent only on non-straight trajectories;
and anchor projections fell from $13.5$ to $0.0$ per episode on shape and from $31.5$ to $1.47$ on
dynamics, so the trust region no longer acts as the controller. The three defects behind the
earlier ``pinned'' allocation --- a normaliser hard-coded to zero at the call site, the
Monte-Carlo epistemic spread used where the predictive variance was fitted, and a raw mean squared
error divided by a temperature fitted on variance-normalised residuals --- are documented in
\S\ref{sec:stls} because each of them silently invalidated the policy while every stage still
reported success.

The success column does not follow. On the shape family STLS-JEPA reaches $38.1\%$ against $39.9\%$
for the frozen model and for DTA-JEPA; on the dynamics family it reaches $18.8\%$ against $17.7\%$
for all three baselines, a difference of one episode out of $96$. This is the falsifiable branch we
pre-registered: \emph{the diagnoses were right and the repairs did not transfer to success}. At this
planning budget the success rate is not set by the allocation policy, the trust region or the
refinement depth --- all of which we can now show behave as designed --- but by something the
allocation cannot reach.

\begin{table}[t]\centering
\caption{STLS-JEPA against the three baselines on identical models and episodes. Diagnostics are the pre-registered observables of Table~\ref{tab:stls}.}
\label{tab:stlsres}\small
\begin{tabular}{llrrrrr}
\toprule
Family & Method & Success (\%) & Steps/upd & Depth & Anchor/ep & Rollback \\
\midrule
shape & Frozen & 39.9 & -- & -- & -- & -- \\
shape & AdaJEPA & 38.1 & 1.00 & -- & -- & 0.00 \\
shape & DTA-JEPA & 39.9 & 2.32 & 2.86 & 13.5 & 0.43 \\
shape & STLS-JEPA & 38.1 & 1.11 & 1.06 & 0.0 & 0.27 \\
\midrule
dyn & Frozen & 17.7 & -- & -- & -- & -- \\
dyn & AdaJEPA & 17.7 & 1.00 & -- & -- & 0.00 \\
dyn & DTA-JEPA & 17.7 & 2.10 & 2.57 & 31.5 & 0.29 \\
dyn & STLS-JEPA & 18.8 & 1.32 & 1.11 & 1.5 & 0.64 \\
\bottomrule
\end{tabular}\end{table}

\subsection{The world model improves; the planner does not follow}

Straightening was introduced to improve the conditioning of the latent cost, and it does change the
world model measurably (Table~\ref{tab:wm}): next-step prediction error falls from $25.80$ to
$19.98$ on push ($-22\%$) and from $46.63$ to $16.64$ on the maze ($-64\%$), the variance head's
ranking and relative calibration both improve on push ($\rho$ $0.549\to0.630$, ECE $0.821\to0.790$),
and the maze depth curve becomes monotonically decreasing ($30.1\to21.3\to16.4\to16.6$) where it
previously saturated at depth three ($51.9\to44.8\to43.3\to46.6$).

A first reading of the closed-loop numbers suggested something stronger: the frozen model's
dynamics-family success on those checkpoints is $17.7\%$ against $35.4\%$ on the originals, a
halving of planning success alongside a $64\%$ error reduction. \emph{We could not reproduce that
under control, and we therefore withdraw the causal reading.} The two checkpoints differ in three
ways at once (straightening weight, learning rate $5$ vs $6\times10^{-4}$, and nothing else), and on
the push family both axes moved the same way (success $32.1\to39.9\%$ with error $-22\%$), which
already argued against a single mechanism.

\subsection{Controlled sweep: the two axes are separate}

Four maze world models were then trained with identical learning rate, epochs and seed, differing
only in $\lambda_{\mathrm{str}}$ (Table~\ref{tab:geom}). Prediction error falls monotonically, by
$72\%$ end to end, and the closed-loop result is unequivocal in the opposite direction: success on
the \emph{unshifted} dynamics is identical at every weight ($15/24$ episodes at each of
$\lambda_{\mathrm{str}}\in\{0,0.02,0.1,0.5\}$, Fisher exact $p=1.000$), and success under the mass
shift is flat up to $\lambda_{\mathrm{str}}=0.1$ ($9/24$ vs $9/24$, $p=1.000$) before falling at
$0.5$ ($9/24\to4/24$, $p=0.193$, not significant at this sample size). A $72\%$ improvement in
next-step accuracy therefore buys nothing measurable in planning success, and costs nothing on the
condition the model was trained for.

We also tested the geometric explanation proposed in the previous subsection and failed to confirm
it. None of the four planning-free geometry metrics tracks the movement: the effective dimension of
the latent drops from $15.5$ to about $5$ at the \emph{first} non-zero weight and then stays flat
while success does not move, curvature is non-monotone in $\lambda_{\mathrm{str}}$
($0.49\to1.30\to1.14\to0.20$), and action sensitivity $\|\partial z/\partial u\|$ \emph{rises}
monotonically with the weight that is supposed to be squeezing it
($1.00\to1.29\times10^{-3}$). The honest status is therefore: prediction accuracy and planning
success are decoupled on this benchmark, and the latent property that gates planning has not been
identified by the candidates we measured.

\begin{table}[t]\centering
\caption{Straightening sweep: identical lr, epochs and seed, only $\lambda_{\mathrm{str}}$ varies. Prediction error falls monotonically while the latent's effective dimension on the maze distribution and the closed-loop success at a mass shift follow it down.}
\label{tab:geom}\small
\begin{tabular}{rrrrrrr}
\toprule
$\lambda_{\mathrm{str}}$ & pred.\ err. & curvature & sep/curv & eff.\ dim & succ.\ default & succ.\ lowMass \\
\midrule
0.00 & 37.96 & 0.494 & 0.322 & 15.54 & 62.5 & 37.5 \\
0.02 & 22.09 & 1.301 & 0.863 & 5.43 & 62.5 & 33.3 \\
0.10 & 23.25 & 1.141 & 0.652 & 4.27 & 62.5 & 37.5 \\
0.50 & 10.67 & 0.203 & 1.270 & 5.00 & 62.5 & 16.7 \\
\bottomrule
\end{tabular}\end{table}

\begin{table}[t]\centering
\caption{World model, planning-free: in-distribution next-step prediction on the evaluation
distributions (192 windows). Straightening improves the model on both families; the closed-loop
comparison across the same two pairs is confounded (see text), which is what motivated the
controlled sweep of Table~\ref{tab:geom}.}
\label{tab:wm}\small
\begin{tabular}{llrrrr}
\toprule
Checkpoint & Family & Error & Pred.\ var. & Spearman $\rho$ & rel.\ ECE \\
\midrule
\texttt{push4}  & push & 25.80 & 22.81 & 0.549 & 0.821 \\
\texttt{push4s} & push & 19.98 & 18.66 & 0.630 & 0.790 \\
\texttt{maze25}  & maze & 46.63 & 32.85 & 0.791 & 0.893 \\
\texttt{maze25s} & maze & 16.64 & 12.48 & 0.706 & 0.796 \\
\bottomrule
\end{tabular}
\end{table}
